\section{Marcos aproximados y pendientes medias}
\label{sec:marcos}

\fuente{A, §5, Lema~5.1, Proposición~5.2 y Teorema~5.3,
pp.~17--20; B, §§4.3--4.4, Lema~4.2 y Proposición~4.3,
pp.~39--44. Reconstruimos la desigualdad lineal y su integración.
La existencia de la filtración de Harder--Narasimhan se utiliza como
resultado clásico citado.}

Las aplicaciones integrales de la Nota~\ref{sec:proyecciones} parten de
retículos pequeños y llegan a un mismo retículo receptor. Para comparar
sus pendientes medias interesa saber cuántas direcciones del receptor
siguen cubiertas cuando se retienen solo ciertas partes de las fuentes.
Una familia muy numerosa puede concentrarse en una sola dirección.
La información pertinente será una estimación colectiva de trazas.

El argumento se desarrolla en dos niveles. Primero se estima la
codimensión del espacio generado por imágenes de subespacios. Después
se eligen esos subespacios mediante umbrales de pendiente, se aplica la
estimación para cada umbral y se integra el déficit resultante. Las
condiciones que hacen posible aplicar este mecanismo a los periodos
automorfos quedan identificadas al final.

\subsection{Operadores de Gram y distribución de las imágenes}

Sean $V_i$ y $H$ espacios de Hilbert complejos de dimensión finita, y
$B_i:V_i\to H$ una familia finita de aplicaciones lineales. Escribimos
\[
 r_i=\dim V_i,\qquad R=\sum_i r_i>0,\qquad D=\dim H>0.
\]
El adjunto $B_i^*$ se caracteriza por la identidad del producto hermitiano.
Para un operador $A$, su norma de Hilbert--Schmidt satisface
\[
 \|A\|_{\mathrm{HS}}^2=\operatorname{Tr}(A^*A).
\]
En bases ortonormales es la suma de los cuadrados de los valores absolutos
de las entradas. La desigualdad de Cauchy--Schwarz para este producto
escalar da
$|\operatorname{Tr}(A^*C)|\le\|A\|_{\mathrm{HS}}\|C\|_{\mathrm{HS}}$.
Un proyector ortogonal de rango $k$ tiene norma de Hilbert--Schmidt
$\sqrt{k}$.

Elegimos números $c_i>0$ y definimos
\begin{equation}
 A_i=c_iB_i^*B_i,\qquad \Delta_i=A_i-I_{V_i},\qquad
 Q=\sum_i c_iB_iB_i^*.
 \label{eq:marcos-operadores}
\end{equation}
Los operadores $A_i$ y $Q$ son autoadjuntos semidefinidos positivos;
$\Delta_i$ puede tener autovalores de ambos signos. La identidad
de trazas para aplicaciones entre espacios de dimensiones distintas da
\[
 T:=\operatorname{Tr}Q=\sum_i\operatorname{Tr}A_i.
\]

Las dos hipótesis del marco son
\begin{align}
 \sum_i\|\Delta_i\|_{\mathrm{HS}}^2&\le eR,
       &&0\le e\le\tfrac14,
       \label{eq:marcos-error-fuentes}\\
 \operatorname{Tr}(Q^2)&\le(1+\eta)\frac{T^2}{D},
       &&\eta\ge0.
       \label{eq:marcos-error-destino}
\end{align}
La primera compara las métricas de las fuentes con las inducidas por
los mapas ponderados. La segunda limita la dispersión de los autovalores
del operador colectivo. En efecto,
\begin{equation}
 \left\|Q-\frac TD I_H\right\|_{\mathrm{HS}}^2
 =\operatorname{Tr}(Q^2)-\frac{T^2}{D}
 \le\eta\frac{T^2}{D}.
 \label{eq:marcos-varianza}
\end{equation}
Si $e=\eta=0$, cada $A_i$ es la identidad y $Q=(R/D)I_H$.
La escala escalar del marco es $R/D$; puede ser distinta de $1$.

\subsection{La desigualdad de codimensión}

\begin{proposicion}[Marco aproximado]
Supongamos \eqref{eq:marcos-error-fuentes} y
\eqref{eq:marcos-error-destino}. Si $N_i\subseteq V_i$ son
subespacios y $\varepsilon\ge0$ satisface
$\sum_i\operatorname{codim}_{V_i}N_i\le\varepsilon R$, entonces
\begin{equation}
 \operatorname{codim}_H\left(\sum_i B_iN_i\right)
 \le(\eta+4\varepsilon+4\sqrt e)D.
 \label{eq:marcos-codimension}
\end{equation}
\label{prop:marco-lineal}
\end{proposicion}

\begin{proof}
Como $T-R=\sum_i\operatorname{Tr}\Delta_i$, dos aplicaciones de
Cauchy--Schwarz dan
\begin{align*}
 |T-R|&\le\sum_i\sqrt{r_i}\|\Delta_i\|_{\mathrm{HS}}\\
 &\le\sqrt R\left(\sum_i\|\Delta_i\|_{\mathrm{HS}}^2\right)^{1/2}
 \le R\sqrt e.
\end{align*}
Por tanto,
\begin{equation}
 T\ge(1-\sqrt e)R\ge R/2>0.
 \label{eq:marcos-traza-positiva}
\end{equation}

Sea $S=\sum_iB_iN_i$. Denotemos por $P$ el proyector ortogonal
sobre $S^\perp$ y por $P_i$ el proyector sobre $N_i^\perp$.
La aplicación $PB_i$ se anula en $N_i$, de modo que $PB_i=PB_iP_i$.
Una proyección ortogonal no aumenta las normas de los vectores; aplicando
ese hecho a una base ortonormal de $V_i$ obtenemos
\begin{align}
 \operatorname{Tr}(PQ)
 &=\sum_i c_i\|PB_i\|_{\mathrm{HS}}^2\notag\\
 &\le\sum_i c_i\|B_iP_i\|_{\mathrm{HS}}^2
 =\sum_i\operatorname{Tr}(P_iA_i)\notag\\
 &=\sum_i\operatorname{rk}P_i+
       \sum_i\operatorname{Tr}(P_i\Delta_i).
 \label{eq:marcos-traza-omitida}
\end{align}
Las trazas del primer miembro y de $P_iA_i$ son reales y no negativas:
por ciclicidad, coinciden con las de $PQP$ y $P_iA_iP_i$.
Esto no exige que los productos $PQ$ o $P_iA_i$ sean autoadjuntos.

Para el término de error,
\begin{align*}
 \left|\sum_i\operatorname{Tr}(P_i\Delta_i)\right|
 &\le\sum_i\sqrt{\operatorname{rk}P_i}\|\Delta_i\|_{\mathrm{HS}}\\
 &\le\left(\sum_i\operatorname{rk}P_i\right)^{1/2}
       \left(\sum_i\|\Delta_i\|_{\mathrm{HS}}^2\right)^{1/2}
 \le R\sqrt e,
\end{align*}
pues $\sum_i\operatorname{rk}P_i\le R$. La hipótesis de codimensión
controla el término de identidad de \eqref{eq:marcos-traza-omitida}, y da
\begin{equation}
 0\le\operatorname{Tr}(PQ)\le(\varepsilon+\sqrt e)R.
 \label{eq:marcos-control-omitido}
\end{equation}

Ponemos $z=\operatorname{Tr}(PQ)/T$. La positividad de las trazas
sobre $S$ y $S^\perp$ implica $0\le z\le1$;
\eqref{eq:marcos-traza-positiva} y \eqref{eq:marcos-control-omitido}
implican $z\le2(\varepsilon+\sqrt e)$. Aplicando Cauchy--Schwarz
a $I_H-P$ y $Q$,
\[
 (1-z)^2T^2
 =\bigl(\operatorname{Tr}((I_H-P)Q)\bigr)^2
 \le\dim S\,\operatorname{Tr}(Q^2)
 \le\dim S\,(1+\eta)T^2/D.
\]
Podemos cancelar $T^2>0$. Además,
\[
 \frac{(1-z)^2}{1+\eta}\ge1-\eta-2z,
\]
porque al multiplicar por $1+\eta$ la diferencia de ambos miembros
es $z^2+\eta^2+2\eta z\ge0$. Por consiguiente,
\[
 \frac{\dim S}{D}\ge1-\eta-4\varepsilon-4\sqrt e,
\]
que equivale a \eqref{eq:marcos-codimension}.
\end{proof}

Esta es la prueba de A, Proposición~5.2, desarrollada con los operadores
de compresión explícitos. B, Lema~4.2, conserva primero la cota más fina
$R(\varepsilon+\sqrt{\varepsilon e})$ y compara la traza omitida
con \eqref{eq:marcos-varianza}; llega a la misma constante.
Para $\varepsilon\ge1$ la conclusión es inmediata de
$\operatorname{codim}S\le D$. Los mapas pueden tener núcleo y algunos
$N_i$ pueden ser cero. Cuando el coeficiente de la derecha es mayor
o igual que $1$, la desigualdad sigue siendo válida, aunque no mejora
esa cota elemental.

\subsection{Ejemplos de concentración y redundancia}

Tomemos $V_i=\C$, $H=\C^D$, $c_i=1$ y $B_i(z)=ze_1$ para
$1\le i\le n$, donde $e_1$ pertenece a la base ortonormal habitual.
Entonces $A_i=1$, así que $e=0$ es admisible, pero
\[
 Q=nP_{\C e_1},\qquad T=n,\qquad
 \operatorname{Tr}(Q^2)=n^2.
\]
La segunda hipótesis exige $\eta\ge D-1$. Aumentar $n$ conserva
esa exigencia: todas las imágenes siguen en la misma recta.

En cambio, si se toman dos copias de cada aplicación
$\C\to\C^D$, $z\mapsto ze_j$, entonces $R=2D$ y $Q=2I_H$.
Se tiene $e=\eta=0$. Hay redundancia: la aplicación agregada
$\bigoplus_iV_i\to H$, $(v_i)_i\mapsto\sum_iB_iv_i$, tiene
núcleo de dimensión $D$. Ese núcleo no es el núcleo de un mapa
individual y es compatible con un marco exacto.

Con una sola copia de cada coordenada y $D=10$, suprimir una fuente
de dimensión uno da $\varepsilon=1/10$. La imagen retenida tiene
codimensión $1$, mientras \eqref{eq:marcos-codimension} garantiza
codimensión a lo sumo $4$. El ejemplo muestra el carácter uniforme
de la constante, que no pretende ser óptima en cada familia.

\subsection{Normalización aritmética de las pendientes}

Sea $K$ un cuerpo de números, de grado $d_K$, y sea $E$ un
retículo hermitiano: un $\OK$-módulo proyectivo de rango finito positivo con
métricas compatibles en sus realizaciones arquimedianas. Para una
sección racional no nula $v$ de $\det E$, el grado normalizado es
\begin{equation}
 \widehat{\deg}_{\mathrm{Ar}}(E)=\frac1{d_K}
 \left(\sum_{\mathfrak p}\operatorname{ord}_{\mathfrak p}(v;\det E)
                 \log N\mathfrak p
       -\sum_{\sigma:K\hookrightarrow\C}\log\|v\|_\sigma\right).
 \label{eq:marcos-grado}
\end{equation}
La primera suma recorre los primos finitos. El orden se mide respecto
de un generador local y puede ser negativo. La fórmula del producto
hace que el resultado sea independiente de $v$; no se necesita una
base integral global de $E$.

Usamos la normalización de A y B:
\begin{equation}
 \mu(E)=\frac{2\widehat{\deg}_{\mathrm{Ar}}(E)}{\operatorname{rk}E},
 \qquad
 h(B)=\frac2{d_K}\sum_{\sigma:K\hookrightarrow\C}
                         \log\|B_\sigma\|_{\mathrm{op}}\quad(B\ne0).
 \label{eq:marcos-normalizacion}
\end{equation}
La suma comprende todas las inmersiones, contando las dos de cada par
conjugado. Para $K$ totalmente complejo es equivalente sumar
$\log\|B_v\|_{\mathrm{op}}^2$ una vez por lugar complejo y multiplicar
por $2/d_K$, como en B. El número $h(B)$ puede ser negativo.
Sobre $\Q$, $\mu(E)=-2\log\operatorname{covol}(E)/\operatorname{rk}E$,
el doble de la pendiente introductoria.

Si $K$ es totalmente complejo, $r_2=d_K/2$ y $r=\operatorname{rk}E$,
la relación con el covolumen euclídeo de Minkowski es
\begin{equation}
 \mathcal H(E)=2^{r_2r}|\Delta_K|^{-r/2}
                       \operatorname{covol}_{\R}(E),
 \qquad \mu(E)=-\frac2{d_Kr}\log\mathcal H(E).
 \label{eq:marcos-covolumen}
\end{equation}
El factor finito normaliza el covolumen de $\OK$; con estas
convenciones $\mu(\OK^r)=0$ para las métricas habituales.
Las fórmulas \eqref{eq:marcos-grado}--\eqref{eq:marcos-covolumen}
son las de A, (5.1)--(5.3), pp.~17--18, y B, p.~39.

Un subretículo $M\subset E$ es \emph{saturado} cuando $E/M$ no
tiene torsión. La saturación reemplaza $M$ por
$E\cap(M\otimes_{\OK}K)$. Si $E'\subset E$ tiene el mismo rango
$r>0$, las mismas métricas y un cociente finito de cardinal $J$, entonces
\begin{equation}
 \widehat{\deg}_{\mathrm{Ar}}(E')=
 \widehat{\deg}_{\mathrm{Ar}}(E)-\frac{\log J}{d_K},
 \qquad \mu(E')=\mu(E)-\frac{2\log J}{d_Kr}.
 \label{eq:marcos-coste-indice}
\end{equation}
Ampliar el retículo invierte el signo de esa diferencia. Esta es la
contribución que conecta los denominadores de la
Nota~\ref{sec:proyecciones} con las pendientes.

\subsection{La filtración de Harder--Narasimhan}

La pendiente máxima $\mu_{\max}(E)$ es el máximo de las pendientes
de los subretículos saturados no nulos. La mínima $\mu_{\min}(E)$
es el mínimo de las pendientes de los cocientes no nulos sin torsión,
provistos de las métricas cociente. Un retículo es \emph{semiestable}
cuando $\mu_{\max}(E)=\mu(E)$.

La teoría clásica proporciona una filtración canónica por subretículos
saturados
\[
 0=E_0\subset E_1\subset\cdots\subset E_\ell=E,
\]
cuyos cocientes sucesivos, con las métricas inducidas, son semiestables
y tienen pendientes estrictamente decrecientes. Repetir cada pendiente
según el rango de su cociente produce una lista
\begin{equation}
 \lambda_1(E)\ge\cdots\ge\lambda_r(E),\qquad
 \sum_{j=1}^r\lambda_j(E)=r\mu(E).
 \label{eq:marcos-lista-hn}
\end{equation}
La primera y la última son $\mu_{\max}(E)$ y $\mu_{\min}(E)$.
La existencia y las propiedades de esta filtración se reciben de la
teoría de reducción aritmética; véanse Grayson \cite{grayson-hn},
la discusión de A, p.~18, y B, p.~42. La aditividad del grado con
métricas de subespacio y cociente está recogida también en
\cite[§4.1.1, Lema~4.2]{bost-pendientes}.

Escribimos $E^{\ge t}$ para la etapa de la filtración que retiene
los cocientes de pendiente al menos $t$. Entonces
\begin{align}
 \operatorname{rk}E^{\ge t}&=\#\{j:\lambda_j(E)\ge t\},
                  \label{eq:marcos-corte-hn}\\
 \mu_{\min}(E^{\ge t})&\ge t,\qquad
 \mu_{\max}(E/E^{\ge t})<t,
                  \label{eq:marcos-extremos-hn}
\end{align}
donde cada pendiente se considera solo si el retículo correspondiente
es no nulo. La desigualdad estricta del cociente resulta de haber
retenido también los cocientes de pendiente igual a $t$.
No se elige una descomposición de $E$ en los cocientes sucesivos:
la filtración y sus rangos bastan para el argumento.

\subsection{Una aplicación integral desplaza los umbrales}

\begin{lema}[Comparación de pendientes]
Si $B:E\to F$ es una aplicación integral no nula entre retículos
hermitianos, entonces
\begin{equation}
 \mu_{\min}(E)\le\mu_{\max}(F)+h(B).
 \label{eq:marcos-morfismo-pendientes}
\end{equation}
Para todo $t\in\R$ se tiene además
\begin{equation}
 B(E^{\ge t})\subseteq F^{\ge t-h(B)}.
 \label{eq:marcos-inclusion-filtrada}
\end{equation}
\label{lem:marcos-altura}
\end{lema}

\begin{proof}
Sea $I=B(E)$ con la métrica inducida por $F$, y sea $q=\operatorname{rk}I$.
El núcleo es saturado porque $E/\ker B\simeq I$ no tiene torsión.
La aplicación inducida $\overline B:E/\ker B\to I$ es un isomorfismo
de módulos. Con la métrica cociente en el dominio, su determinante
tiene norma a lo sumo $\|B_\sigma\|_{\mathrm{op}}^q$ en cada
inmersión. La fórmula del grado, aplicada a ese determinante, da
\[
 \mu(E/\ker B)\le\mu(I)+h(B).
\]
La saturación $I^{\mathrm{sat}}$ en $F$ tiene el mismo rango y grado
mayor o igual. Por las definiciones de pendiente mínima y máxima,
\[
 \mu_{\min}(E)\le\mu(E/\ker B)
 \le\mu(I)+h(B)\le\mu(I^{\mathrm{sat}})+h(B)
 \le\mu_{\max}(F)+h(B).
\]

Para probar la inclusión, compongamos la restricción de $B$ a
$E^{\ge t}$ con la aplicación cociente
$F\to F/F^{\ge t-h(B)}$. Si la composición fuera no nula, su
altura sería a lo sumo $h(B)$, porque ni restringir el dominio con
su métrica inducida ni proyectar con la métrica cociente aumenta
la norma de operador. Las desigualdades
\eqref{eq:marcos-morfismo-pendientes} y
\eqref{eq:marcos-extremos-hn} darían
\[
 t\le\mu_{\min}(E^{\ge t})
 \le\mu_{\max}(F/F^{\ge t-h(B)})+h(B)<t,
\]
una contradicción. La composición es cero, lo que prueba
\eqref{eq:marcos-inclusion-filtrada}.
\end{proof}

Este es el Lema~5.1 de A. La comparación mediante determinantes es
una forma del método de pendientes de Bost
\cite[§4.1.4, pp.~199--200]{bost-pendientes}. Se han mantenido
las métricas cociente y la saturación para que la contribución integral
y la norma arquimediana correspondan a la misma aplicación.

La integridad es una hipótesis efectiva. En rango uno sobre $\Q$, con
retículos y métricas habituales, multiplicar por $1/p$ tiene altura
arquimediana $-2\log p$, pero ambas pendientes son cero. Aplicar
\eqref{eq:marcos-morfismo-pendientes} a ese mapa racional daría
$0\le-2\log p$, que es falso. El denominador introduce un coste
finito que la hipótesis integral evita omitir.

\subsection{La proporción de pendientes bajas en las fuentes}

Sean $E_i$ retículos de rangos $r_i>0$, $R=\sum_i r_i$ y
$h_i\in\R$. La lista de pendientes desplazadas es
\[
 a_{i,j}=\lambda_j(E_i)-h_i\qquad(1\le j\le r_i).
\]
Cada término recibe peso $1/R$. La media de esta lista es
\begin{equation}
 \frac1R\sum_{i,j}a_{i,j}
 =\sum_i\frac{r_i}{R}\bigl(\mu(E_i)-h_i\bigr).
 \label{eq:marcos-media-fuentes}
\end{equation}
Los pesos $r_i/R$ proceden de contar dimensiones. Sustituirlos por
probabilidades arbitrarias modificaría la relación con la hipótesis
de marco, que está normalizada por $R$.

\begin{lema}[Cola de una lista de pendientes]
Supongamos que la media de los $a_{i,j}$ es al menos $s$ y que
todos satisfacen $a_{i,j}\le s+u$, con $u\ge0$. Para $\tau>0$, sea
\[
 q(\tau)=\frac1R\#\{(i,j):a_{i,j}<s-\tau\}.
\]
Entonces
\begin{equation}
 q(\tau)\le\frac{u}{u+\tau}.
 \label{eq:marcos-cola-fuentes}
\end{equation}
Si $u=0$, todos los términos son $s$ y $q(\tau)=0$.
\label{lem:marcos-cola}
\end{lema}

\begin{proof}
La parte de la lista contada por $q(\tau)$ está por debajo de
$s-\tau$, y el resto está a lo sumo en $s+u$. Por tanto,
\[
 s\le\frac1R\sum_{i,j}a_{i,j}
 \le q(\tau)(s-\tau)+(1-q(\tau))(s+u)
 =s+u-q(\tau)(u+\tau).
\]
Como $u+\tau>0$, esta desigualdad implica
\eqref{eq:marcos-cola-fuentes}. Si $u=0$, una lista cuyos términos
son a lo sumo $s$ y cuya media es al menos $s$ tiene todos sus
términos iguales a $s$.
\end{proof}

La cota superior individual controla cuánto pueden compensar las
pendientes altas a las bajas. La media por sí sola no limita esa
proporción: una sola pendiente arbitrariamente grande podría compensar
muchas pendientes pequeñas.

\subsection{El teorema de transferencia de la pendiente media}

\begin{teorema}[Transferencia mediante un marco aproximado]
Sea $K$ un cuerpo de números totalmente complejo. Sean $E_i$ una
familia finita no vacía de retículos hermitianos de rangos $r_i>0$,
y $F$ otro de rango $r>0$. Sean $b_i:E_i\to F$ aplicaciones
integrales y $h_i\in\R$ tales que $h(b_i)\le h_i$ cuando $b_i\ne0$.
Escribamos $R=\sum_i r_i$. Supongamos que, para ciertos
$s\in\R$ y $u,H_0\ge0$,
\begin{align}
 \sum_i\frac{r_i}{R}\bigl(\mu(E_i)-h_i\bigr)&\ge s,
                   \label{eq:marcos-hip-media}\\
 \mu_{\max}(E_i)-h_i&\le s+u\quad\text{para todo }i,
                   \label{eq:marcos-hip-max}\\
 \mu_{\min}(F)&\ge s-H_0.
                   \label{eq:marcos-hip-min}
\end{align}
En una inmersión compleja fijada, elijamos estructuras de Hilbert
auxiliares, un espacio $\widetilde H$ de dimensión $D>0$, una
aplicación lineal sobreyectiva $\pi:\widetilde H\to F_\sigma$ y
levantamientos $\widetilde B_i:(E_i)_\sigma\to\widetilde H$ con
$\pi\widetilde B_i=(b_i)_\sigma$.
Supongamos que los levantamientos satisfacen
\eqref{eq:marcos-error-fuentes}--\eqref{eq:marcos-error-destino}
con coeficientes positivos, $0\le e\le1/4$ y $\eta\ge0$, y que
\begin{equation}
 r\ge(1-\delta)D,\qquad 0\le\delta<1.
 \label{eq:marcos-hip-proyeccion}
\end{equation}
Entonces
\begin{equation}
 \boxed{\quad
 \mu(F)\ge s-
 \frac{4u\log(1+H_0/u)+(\eta+4\sqrt e)H_0}{1-\delta}.
 \quad}
 \label{eq:marcos-transferencia}
\end{equation}
Cuando $u=0$, el término $u\log(1+H_0/u)$ se interpreta como $0$.
\label{teo:marcos-transferencia}
\end{teorema}

\begin{proof}
Fijemos $\tau>0$. En cada fuente retenemos la realización compleja
del corte de Harder--Narasimhan
\[
 N_i(\tau)=(E_i^{\ge s-\tau+h_i})_\sigma.
\]
Su codimensión cuenta las pendientes desplazadas menores que
$s-\tau$. Las hipótesis \eqref{eq:marcos-hip-media} y
\eqref{eq:marcos-hip-max}, junto con el Lema~\ref{lem:marcos-cola},
dan
\begin{equation}
 \frac1R\sum_i\operatorname{codim}N_i(\tau)
 =q(\tau)\le\frac{u}{u+\tau}.
 \label{eq:marcos-descarte-fuentes}
\end{equation}

Si $b_i\ne0$, el Lema~\ref{lem:marcos-altura} implica
\[
 b_i(E_i^{\ge s-\tau+h_i})
 \subseteq F^{\ge s-\tau+h_i-h(b_i)}
 \subseteq F^{\ge s-\tau}.
\]
La segunda inclusión usa $h_i\ge h(b_i)$ y que los cortes disminuyen
al aumentar el umbral. Para $b_i=0$ la inclusión es automática.

Aplicamos la Proposición~\ref{prop:marco-lineal} a los levantamientos
y a $N_i(\tau)$. Si $\widetilde S=\sum_i\widetilde B_iN_i(\tau)$,
entonces
\[
 \operatorname{codim}_{\widetilde H}\widetilde S
 \le\left(\eta+4\sqrt e+\frac{4u}{u+\tau}\right)D.
\]
La aplicación sobreyectiva $\pi$ induce una sobreyección
$\widetilde H/\widetilde S\to F_\sigma/\pi\widetilde S$.
Por ello proyectar no aumenta la codimensión. Además,
$\pi\widetilde S\subseteq(F^{\ge s-\tau})_\sigma$.
Al dividir por $r\ge(1-\delta)D>0$ obtenemos
\begin{equation}
 q_F(\tau):=\frac{\#\{j:\lambda_j(F)<s-\tau\}}r
 \le\frac1{1-\delta}
       \left(\eta+4\sqrt e+\frac{4u}{u+\tau}\right).
 \label{eq:marcos-cola-destino}
\end{equation}
Esta parte del razonamiento utiliza las métricas auxiliares solo para
aplicar el resultado lineal. Las alturas $h(b_i)$ se siguen calculando
con las métricas aritméticas originales y con todas las inmersiones.
La sobreyección $\pi$ no necesita ser ortogonal ni tener una norma acotada.

Por \eqref{eq:marcos-hip-min}, cada pendiente del destino satisface
$\lambda_j(F)\ge s-H_0$. Con $a_+=\max(a,0)$,
\begin{equation}
 (s-\lambda_j(F))_+
 =\int_0^{H_0}\mathbf1_{\{\lambda_j(F)<s-\tau\}}\dd\tau.
 \label{eq:marcos-capas}
\end{equation}
La igualdad se comprueba midiendo el intervalo
$0<\tau<s-\lambda_j(F)$ cuando el extremo es positivo.
Las posibles igualdades en los extremos no afectan a la integral.
Al promediar una suma finita,
\begin{equation}
 s-\mu(F)\le\frac1r\sum_j(s-\lambda_j(F))_+
 =\int_0^{H_0}q_F(\tau)\dd\tau.
 \label{eq:marcos-deficit-medio}
\end{equation}
Si $u>0$, integrar \eqref{eq:marcos-cola-destino} y usar
\[
 \int_0^{H_0}\frac{u}{u+\tau}\dd\tau
 =u\log(1+H_0/u)
\]
prueba \eqref{eq:marcos-transferencia}.
Si $u=0$, $q(\tau)=0$ para todo $\tau>0$ y la misma integración
solo conserva el término $(\eta+4\sqrt e)H_0/(1-\delta)$.
Si $H_0=0$, todas las pendientes de $F$ son al menos $s$ y la
conclusión se sigue directamente.
\end{proof}

La fórmula coincide con A, Teorema~5.3, y B, Proposición~4.3.
El número de mapas no aparece en la pérdida. La información sobre
toda la familia se reúne mediante los rangos y los dos defectos del
marco; no se introduce un término $\log|I|$ al integrar.

\subsection{Cómo la integral recupera un déficit}

Consideremos la lista de pendientes $(-2,0,1,3)$, con $s=1$ y
$H_0=3$. La proporción por debajo de $1-\tau$ y su contribución
a la integral son
\begin{center}
\begin{tabular}{@{}ccc@{}}
\toprule
Intervalo & $q_F(\tau)$ & Área \\
\midrule
$0<\tau<1$ & $1/2$ & $1/2$ \\
$1\le\tau<3$ & $1/4$ & $1/2$ \\
$\tau\ge3$ & $0$ & $0$ \\
\bottomrule
\end{tabular}
\end{center}
El área total es $1$, igual al déficit positivo medio
$(3+1+0+0)/4$. La media de la lista es $1/2$, de modo que
$s-\mu=1/2$. La diferencia entre ambas cantidades es el excedente
positivo medio $(0+0+0+2)/4=1/2$. Por eso
\eqref{eq:marcos-deficit-medio} comienza con una desigualdad.
El ejemplo ilustra la identidad de integración; por sí solo no
especifica mapas que satisfagan las hipótesis del teorema.

\subsection{El coste logarítmico y el tamaño de los errores}

\begin{lema}
Si $u,U,H_0,b,X\ge0$, $u\le U$ y $H_0\le Xb$, entonces
\begin{equation}
 u\log(1+H_0/u)\le(U+b)\log(1+X),
 \label{eq:marcos-absorcion-log}
\end{equation}
con la misma convención cuando $u=0$.
\end{lema}

\begin{proof}
Para $u>0$, la recta tangente a $\log y$ en $y=1+X$ da
\[
 \log(1+H_0/u)
 \le\log(1+X)+\frac{H_0/u-X}{1+X}.
\]
Multiplicar por $u$, usar $H_0\le Xb$ y descartar el término
no positivo $-uX/(1+X)$ produce
\[
 u\log(1+H_0/u)
 \le u\log(1+X)+\frac{Xb}{1+X}.
\]
La desigualdad $X/(1+X)\le\log(1+X)$ se deduce de que
la diferencia vale cero en $X=0$ y tiene derivada
$X/(1+X)^2\ge0$. Con $u\le U$ obtenemos
\eqref{eq:marcos-absorcion-log}. Para $u=0$ el miembro izquierdo
es cero y el derecho es no negativo.
\end{proof}

Este es el Lema~5.4 de A, p.~20. Si $X$ es una potencia fija de
$x$, su contribución tiene un factor $\log x$. Sustituirla
prematuramente por un producto $uH_0$ perdería esa información.
Los errores del marco aparecen multiplicados por $H_0$, mientras
que el descarte por pendientes se integra mediante la función
$u/(u+\tau)$. La pérdida $1/(1-\delta)$ procede de comparar
la dimensión auxiliar con el rango del receptor.

\subsection{Qué debe aportar la construcción de periodos}

La Nota~\ref{sec:proyecciones} construye formalmente aplicaciones
integrales $b_\alpha:N_\alpha\to C$ a partir de la divisibilidad
de los emparejamientos. El Teorema~\ref{teo:marcos-transferencia}
indica qué información adicional debe acreditarse para utilizarlas:

\begin{center}
\small
\begin{tabularx}{\textwidth}{@{}>{\raggedright\arraybackslash}p{0.32\textwidth}Y@{}}
\toprule
Dato de entrada & Función en el argumento \\
\midrule
Integridad de $b_\alpha$ y altura $h_\alpha$ &
Transportar el corte de pendiente sin omitir índices ni denominadores. \\
Media de $\mu(N_\alpha)-h_\alpha$ y cota de $\mu_{\max}(N_\alpha)-h_\alpha$ &
Fijar $s,u$ y controlar la fracción descartada en cada umbral. \\
Mínimo de las pendientes de $C$ &
Fijar el extremo $H_0$ de la integral del déficit. \\
Defectos $e,\eta$ del marco &
Controlar la pérdida de direcciones tras restringir simultáneamente las fuentes. \\
Rango retenido y dimensión auxiliar &
Fijar $\delta$ y comparar codimensión con proporción de pendientes del receptor. \\
\bottomrule
\end{tabularx}
\end{center}

Estas cotas tienen que referirse a los mismos mapas y a normalizaciones
compatibles. Las normas individuales no proporcionan por sí solas
la estimación colectiva de $\operatorname{Tr}(Q^2)$, como muestra
el ejemplo de concentración. Un descarte previo de fuentes que cambie
$e$ o $\eta$ tampoco se identifica con el parámetro $\delta$, que
mide una pérdida de dimensión en el destino. B, pp.~44--46, distingue
esos pasos en la aplicación a periodos.

La desigualdad abstracta queda demostrada con las entradas indicadas.
Producirlas para la familia automorfa requiere estudiar las torres theta,
las estimaciones de evaluación, las estructuras integrales y los costes
de las proyecciones de A, §§6--8, y B, §§4.1--4.5 y~5.
Esta nota no da todavía la cota del exceso geométrico de la
Nota~\ref{sec:caracteres} ni la ventana aritmética de la
Nota~\ref{sec:ventanas}.

\subsection{Concordancia localizada con Lean}
\label{subsec:marcos-lean}

El cotejo usa el espejo citado en \cite{lean}, fijado en la revisión
\href{https://github.com/CocoSgt/riemann-hypothesis-lean/tree/7cbf7a0c32fca1badf05931879a37c3f8a2b03b4}{\texttt{7cbf7a0}},
cuyo identificador completo figura en el registro.
Se leyeron catorce archivos de los módulos de Harder--Narasimhan y
de corangos. El registro conserva sus rutas, declaraciones, hipótesis,
URL fijadas y hashes. Los nombres siguientes llevan el prefijo
\texttt{Formalism.HN}, salvo que se indique otro.

\begin{enumerate}
\item \textbf{La constante del marco lineal.}
En
\href{https://github.com/CocoSgt/riemann-hypothesis-lean/blob/7cbf7a0c32fca1badf05931879a37c3f8a2b03b4/src/Formalism/HarderNarasimhan/StatementFrameF.lean}{\texttt{StatementFrameF.lean}},
\texttt{FrameLA1} define el enunciado
\eqref{eq:marcos-codimension}. Su prueba es
\href{https://github.com/CocoSgt/riemann-hypothesis-lean/blob/7cbf7a0c32fca1badf05931879a37c3f8a2b03b4/src/Formalism/HarderNarasimhan/HarderNarasimhanProofs/FrameLinearAlgebraFirst.lean}{\texttt{frame\_linear\_algebra\_first}}.
Recibe los coeficientes positivos, $R>0$, las dos hipótesis de trazas
y la cota de codimensión total. Expresa la conclusión como la dimensión
del complemento ortogonal. La no negatividad de $\varepsilon$ se deduce
de las hipótesis; el caso $D=0$ se trata aparte. La prueba formal leída
sigue el cálculo de varianza de B. La prueba por $I-P$ desarrollada
arriba sigue A; las conclusiones y las constantes coinciden.

\item \textbf{La transferencia precisa.}
En el mismo archivo, \texttt{TheoremFSharp} define el enunciado
preciso, cuya prueba es
\href{https://github.com/CocoSgt/riemann-hypothesis-lean/blob/7cbf7a0c32fca1badf05931879a37c3f8a2b03b4/src/Formalism/HarderNarasimhan/HarderNarasimhanProofs/TheoremFSharp.lean}{\texttt{frame\_theorem\_sharp}}.
La interfaz pública
\href{https://github.com/CocoSgt/riemann-hypothesis-lean/blob/7cbf7a0c32fca1badf05931879a37c3f8a2b03b4/src/Formalism/HarderNarasimhan/ShiftInvariance.lean}{\texttt{theoremFSharp}}
la expone directamente. Sus hipótesis incluyen la integridad, las alturas
de los mapas no nulos, la media ponderada por $r_i$, los extremos de
pendiente, las dos cotas de marco y la relación de rangos.
El levantamiento se escribe por matrices: su bloque superior representa
el mapa aritmético después de cambios de coordenadas invertibles.
La proyección sobre ese bloque realiza la sobreyección de la prueba.
La interfaz admite fuentes individuales de rango cero si $R>0$;
nuestra exposición usa el caso de rangos positivos de A y B.

\item \textbf{La normalización y los subobjetos.}
\href{https://github.com/CocoSgt/riemann-hypothesis-lean/blob/7cbf7a0c32fca1badf05931879a37c3f8a2b03b4/src/Formalism/HarderNarasimhan/StatementPl.lean}{\texttt{StatementPl.lean}}
define \texttt{ht}, \texttt{size}, \texttt{muDef} y \texttt{htMap}.
La identidad demostrada por \texttt{size\_maximal\_slope} identifica
\texttt{muDef} con \eqref{eq:marcos-covolumen}. Internamente,
si $\lambda$ es una pendiente calculada mediante el covolumen euclídeo,
la pendiente de nuestra normalización es $2\lambda+c_K$, donde
\[
 c_K=-\log2+\frac{\log|\Delta_K|}{d_K}.
\]
Este término ya está incorporado en \texttt{muDef}.
Los extremos y los cortes recorren submódulos sobre $\OK$.
La realización sobre $\Z$ sirve para calcular rango y covolumen;
no amplía la familia de subobjetos a todos los subretículos sobre $\Z$.
Las multiplicidades de las listas formales usan
$\operatorname{rk}_{\Z}=d_K\operatorname{rk}_{\OK}$; el factor común
$d_K$ cancela en las medias y en las proporciones.

\item \textbf{Los cortes y su integral.}
En la prueba de transferencia, \texttt{thmF\_share} estima la
fracción de pendientes bajas y \texttt{thmF\_place} aplica el marco
a los cortes y compara los corangos tras proyectar.
El paso de altura a pendiente de la imagen utiliza
\texttt{slope\_inequality\_mean}; su prueba leída cambia métricas
para igualar las normas en los lugares y aplica una desigualdad
euclídea importada. La reconstrucción anterior usa determinantes
y saturación. Finalmente,
\href{https://github.com/CocoSgt/riemann-hypothesis-lean/blob/7cbf7a0c32fca1badf05931879a37c3f8a2b03b4/src/Formalism/CorankShare/CorankShareProofs/LayerCakeFormula.lean}{\texttt{layer\_cake\_formula}},
del módulo de corangos, integra los recuentos
estrictos, y \texttt{frameB\_integral} calcula el término logarítmico.
La prueba trata expresamente $u=0$ y los mapas nulos.
\end{enumerate}

Hay variantes que deben distinguirse: \texttt{FrameLA}, en el módulo
de corangos, tiene coeficiente $2\eta$; \texttt{TheoremF} resta además
$\log|I|$. Las declaraciones precisas para las fórmulas de esta nota
son \texttt{FrameLA1} y \texttt{TheoremFSharp}.

En este tramo no se encontró una discrepancia de signos o constantes
con A y B. Se distinguieron las definiciones de tipo \texttt{Prop}
de las pruebas leídas. El cotejo no desplegó toda la construcción de la
filtración ni todas las dependencias importadas, y no incluyó compilación,
ejecución de Lean o comparación integral con la distribución original.
Las hipótesis que debe producir la familia automorfa siguen siendo las
entradas de la subsección anterior.

\subsection{Ejercicios con soluciones}

\begin{ejercicio}\label{ej:marcos-proyeccion}
Sea $\pi:W\to V$ una aplicación lineal sobreyectiva entre espacios
vectoriales de dimensión finita, y sea $S\subset W$ un subespacio.
Demuestre que
$\operatorname{codim}_V\pi S\le\operatorname{codim}_W S$.
Determine cuándo hay igualdad.
\end{ejercicio}

\paragraph{Solución del Ejercicio~\ref{ej:marcos-proyeccion}.}
La fórmula de dimensiones da
\[
 \operatorname{codim}_V\pi S
 =\operatorname{codim}_W S-\dim\ker\pi+\dim(S\cap\ker\pi).
\]
El término restado es no negativo. Hay igualdad si y solo si
$\ker\pi\subseteq S$. El cálculo usa sobreyectividad y dimensiones,
sin recurrir a ninguna métrica.

\begin{ejercicio}\label{ej:marcos-pesos}
Dos fuentes tienen rangos $1$ y $9$, y pendientes desplazadas
constantes $9$ y $-1$, respectivamente. Calcule la media de la lista
completa y compárela con la media que asigna peso $1/2$ a cada fuente.
\end{ejercicio}

\paragraph{Solución del Ejercicio~\ref{ej:marcos-pesos}.}
La lista consta de un $9$ y nueve valores $-1$, así que su media es
$(9-9)/10=0$. La media de las dos fuentes con pesos iguales sería
$(9-1)/2=4$. El segundo número no es la media que controla la
codimensión total de los cortes: el marco cuenta las diez dimensiones.

\begin{ejercicio}\label{ej:marcos-altura-negativa}
Sobre $\Q$, sea $E=F=\Z$ como módulo, pero con normas
$\|1\|_E=1$ y $\|1\|_F=e^{-a/2}$, donde $a>0$.
Para la identidad integral $B:E\to F$, calcule $h(B)$ y las
pendientes en la normalización duplicada. Compruebe
\eqref{eq:marcos-morfismo-pendientes}.
\end{ejercicio}

\paragraph{Solución del Ejercicio~\ref{ej:marcos-altura-negativa}.}
La norma de operador es $e^{-a/2}$, por lo que $h(B)=-a$.
Se tiene $\mu(E)=0$ y $\mu(F)=a$; como los rangos son uno,
las pendientes máxima, mínima y media coinciden. La desigualdad
es $0\le a-a=0$. Este ejemplo conserva la integridad del mapa
y muestra por qué la altura arquimediana puede ser negativa.
